\documentclass[10pt]{article}
\usepackage[T1]{fontenc}
\usepackage[utf8]{inputenc}
\usepackage[french]{babel}
\usepackage{amsmath,amssymb,mathrsfs,mathpazo}
\usepackage{geometry}
\geometry{a4paper,margin=2cm}
\usepackage{graphicx,array,tabularx,tabularray,colortbl,multirow,multicol,booktabs,enumitem,pifont,ulem,fancybox,diagbox,cancel}
\usepackage{pstricks,pstricks-add,pst-plot,pst-tree,pst-node,pst-text,pst-func,pst-math,pst-bspline,pst-3dplot,pst-eucl,pst-solides3d}
\let\pscalculate\undefined
\usepackage{pst-bezier}
\usepackage{tikz,pgfplots,tkz-tab}
\usetikzlibrary{arrows,arrows.meta,patterns,calc,positioning,trees,decorations.markings,shapes.geometric,shapes.misc,shapes.arrows,intersections,decorations.pathmorphing,decorations.pathreplacing,backgrounds,shadows,babel}
\pgfplotsset{compat=1.18}
\usepgfplotslibrary{fillbetween}
\usepackage[np]{numprint}
\usepackage{listings,makecell,pas-tableur,twoopt}
\let\barre\undefined
\usepackage[pstricks]{tablvar}
\definecolor{darkgreen}{RGB}{0,100,0}
\providecommand{\red}{\color{red}}
\providecommand{\blue}{\color{blue}}
\providecommand{\green}{\color{green}}
\providecommand{\black}{\color{black}}
\providecommand{\white}{\color{white}}
\providecommand{\gray}{\color{gray}}
\providecommand{\cyan}{\color{cyan}}
\providecommand{\magenta}{\color{magenta}}
\providecommand{\yellow}{\color{yellow}}
\providecommand{\violet}{\color{violet}}
\providecommand{\orange}{\color{orange}}
\providecommand{\euro}{\texteuro}
\providecommand{\EUR}{\texteuro}
\providecommand{\ieme}{\textsuperscript{e}}
\providecommand{\degres}{\ensuremath{^\circ}}
\providecommand{\degree}{\ensuremath{^\circ}}
\providecommand{\vv}[1]{\overrightarrow{#1}}
\providecommand{\Coord}[1]{#1}
\providecommand{\decosix}{}
\providecommand{\starredbullet}{\ensuremath{\bullet}}
\pagestyle{empty}
\setlength{\parindent}{0pt}
\usepackage{adjustbox}\DeclareUnicodeCharacter{2236}{:}\DeclareUnicodeCharacter{2190}{\ensuremath{\leftarrow}}
\providecommand{\R}{\mathbb{R}}
\providecommand{\N}{\mathbb{N}}
\providecommand{\Z}{\mathbb{Z}}
\providecommand{\Q}{\mathbb{Q}}
\providecommand{\D}{\mathbb{D}}
\providecommand{\C}{\mathbb{C}}
\providecommand{\e}{{\mathrm{e}}}
\providecommand{\vect}[1]{\overrightarrow{#1}}
\providecommand{\vectt}[1]{\overrightarrow{\mathrm{#1}}}
\providecommand{\vvt}[1]{\overrightarrow{\mathrm{#1}}}
\providecommand{\np}[2][]{{#2~#1}}
\providecommand{\vv}[1]{{\overrightarrow{#1}}}
\renewcommand{\d}{{\mathrm{d}}}
\providecommand{\degres}{{^\circ}}
\providecommand{\degre}{{^\circ}}
\providecommand{\blue}{\color{blue}}
\providecommand{\red}{\color{red}}
\providecommand{\enskip}{\quad}
\providecommand{\llbracket}{\lbrack\!\lbrack}
\providecommand{\rrbracket}{\rbrack\!\rbrack}
\providecommand{\euro}{{\text{€}}}
\providecommand{\no}{{\text{n°}}}
\providecommand{\interval}[2]{{[#1;#2]}}
\providecommand{\intervFO}[2]{{[#1;#2[}}
\providecommand{\dotfill}{\ldots}

\providecommand{\ds}{\displaystyle}
\providecommand{\textding}[1]{#1}
\providecommand{\Oij}{\(\mathrm{O};\vec{\imath},\vec{\jmath}\)}
\providecommand{\Oijk}{\(\mathrm{O};\vec{\imath},\vec{\jmath},\vec{k}\)}
\providecommand{\Ouv}{\(\mathrm{O};\vec{u},\vec{v}\)}
\providecommand{\intervFF}[2]{\ensuremath{[#1;#2]}}
\renewcommand{\euro}{\ifmmode\text{\texteuro}\else\texteuro\fi}
\DeclareUnicodeCharacter{00B0}{\ifmmode^{\circ}\else\textdegree\fi}
\DeclareUnicodeCharacter{20AC}{\ifmmode\text{\texteuro}\else\texteuro\fi}
% Énoncés : Éducation nationale. Transcription et corrigés : APMEP, Denis Vergès et contributeurs. Découpage et classement : CoopMaths, Sébastien Lozano, Éric Elter et contributeurs. Adaptation de la présentation : Maths Explorer.
% CC BY-SA : https://coopmaths.fr/www/about/
% https://coopmaths.fr/alea/static/bac/2025/tex/bac_2025_06_sujet1_metropole_3_cor.tex
\begin{document}
\noindent\textit{APMEP / CoopMaths. CC BY-SA. Présentation adaptée par Maths Explorer.}\par\medskip
\medskip

\begin{enumerate}
	\item \textbf{Affirmation 1 : Vraie.}

	La lecture de la représentation paramétrique donnée indique que la droite représentée :
	\begin{itemize}
		\item passe (pour le paramètre $t=0$) par le point de coordonnées $(3 ~;~ 2 ~;~ -1)$, autrement dit par B;
		\item et qu'elle est dirigée par le vecteur $\vect{u} \begin{pmatrix}-2 \\-1 \\ 3\end{pmatrix}$.

		Comme on a :\quad $\vect{\mathrm{AB}} \begin{pmatrix} 3 - (-1) \\ 2 - 0 \\ -1 - 5\end{pmatrix}$, soit $\vect{\mathrm{AB}} \begin{pmatrix} 4 \\ 2 \\ -6\end{pmatrix}$, on remarque que $\vect{\mathrm{AB}} = -2\vect{u}$, donc la droite dont on a la représentation paramétrique est dirigée par un vecteur colinéaire à $\vect{\mathrm{AB}}$.
	\end{itemize}

	Ainsi, la représentation paramétrique donnée est celle d'une droite passant par B et dirigée par $\vect{\mathrm{AB}}$ : c'est bien la droite (AB).

	\medskip

	\textbf{Affirmation 2 : Fausse.}

	Comme O est l'origine du repère, on a $\vect{\mathrm{OA}}\begin{pmatrix} -1 \\ 0 \\ 5	\end{pmatrix}$ et $\vect{\mathrm{OB}}\begin{pmatrix} 3 \\ 2 \\ -1	\end{pmatrix}$.

	Ces vecteurs sont clairement non colinéaires, et donc O, A et B définissent bien un plan. Le repère est orthonormal, donc on peut calculer des produits scalaires à l'aide des coordonnées :

	On a bien :\quad $\vect{n} \cdot \vect{\mathrm{OA}} = 5 \times (-1) + (-2) \times 0 + 1 \times 5 = -5 + 0 + 5 = 0$ et $\vect{n}$ est donc bien orthogonal à $\vect{\mathrm{OA}}$,

	par contre, on a :\quad $\vect{n} \cdot \vect{\mathrm{OB}} = 5 \times 3 + (-2) \times 2 + 1 \times (-1) = 15 - 4 - 1 = 10 \neq 0$ et $\vect{n}$ n'est donc pas orthogonal à $\vect{\mathrm{OB}}$.

	Ainsi, $\vect{n}$ n'est pas orthogonal au plan (OAB), car il n'est pas orthogonal à deux vecteurs définissant une base de ce plan.

	\item \textbf{Affirmation 3 : Fausse.}

	Par lecture des représentations paramétriques, les deux droites $d$ et $d'$ sont dirigées respectivement par $\vect{v} \begin{pmatrix} 1 \\-1\\2\end{pmatrix}$ et  $\vect{v'} \begin{pmatrix} 4 \\4\\-6\end{pmatrix}$.

	Ces vecteurs sont clairement non colinéaires, donc les droites ne sont pas parallèles (ni confondues, ni strictement parallèles). Elles peuvent donc être soit sécantes, soit non coplanaires.

	Cherchons s'il existe un coupe $(k~;~s)$ tel que le point de paramètre $k$ sur $d$ soit confondu avec le point de paramètre $s$ sur $d'$, pour cela, résolvons :

	\[
	\begin{aligned}
		\begin{cases}
			15+k=1+4 s \\
			8-k=2+4 s \\
			-6+2 k=1-6 s
		\end{cases}
		&\iff \begin{cases}
			k=-14 + 4 s \\
			8-(-14 + 4 s)=2+4 s \\
			-6+2(- 14 + 4 s)=1-6 s
		\end{cases}\\
		&\iff \begin{cases}
			k=-14+4 s \\
			8+14-4 s=2+4 s \\
			-6-28+8 s=1-6 s
		\end{cases}\\
		&\iff \begin{cases}
			k=-14+4 s \\
			8 s=20 \\
			14 s=35
		\end{cases}\\
		&\iff \begin{cases}
			k=-14+4 s \\
			s=\frac{20}{8}=\frac{5}{2} \\
			s=\frac{35}{14}=\frac{5}{2}
		\end{cases}\\
		&\iff \begin{cases}
			s=2,5 \\
			s=2,5 \\
			k=-4
		\end{cases}
	\end{aligned}
	\]

	Le système a une solution, donc les deux droites ont un point commun, le point qui a comme coordonnées $(11 ~;~ 12 ~;~ -14)$, c'est le point de paramètre $k = -4$ sur $d$ et c'est aussi le point de paramètre $s = 2,5$ sur $d'$.

	Elles sont donc sécantes en ce point, et donc sont coplanaires.

	\item \textbf{Affirmation 4 : Vraie.}

	Le repère est orthonormé, donc par lecture de l'équation de plan, $P$ admet comme vecteur normal $\vect{n'}$ de coordonnées : \quad$\vect{n'} \begin{pmatrix} 1\\-1\\1	\end{pmatrix}$.

	Ainsi, la droite $\Delta$, passant par C et orthogonale à $P$ aura pour représentation paramétrique : \quad $\begin{cases} x = 2 + l\\y = -1 - l\\ z= 2 + l\end{cases}\quad$ avec $l\in\R$.

	Si on appelle H le projeté orthogonal de C sur $P$, alors $\mathrm{H} = P \cap \Delta$. Donc H sera le point de $\Delta$ dont le paramètre fait que les coordonnées paramétriques vérifient l'équation de $P$.

	On résout donc :  \quad $\begin{aligned}[t]
		(2 + l) - (-1 - l) + (2 + l) + 1 = 0 &\iff 3l + 6 = 0\\
		&\iff 3l = -6\\
		&\iff l = -2
	\end{aligned}$

	H est donc le point de paramètre $l = -2$ sur $\Delta$, donc ses coordonnées sont :

	$\mathrm{H} (0 ~;~ 1 ~;~ 0)$.

	La distance CH vaut donc : \quad $\begin{aligned}[t]
		\mathrm{CH} &= \sqrt{(0 - 2)^2 + (1 - (-1))^2 + (0 - 2)^2}\\
			&=\sqrt{4 + 4 + 4}\\
			&=\sqrt{4\times 3}\\
			&=\sqrt{4}\times \sqrt{3}\\
			&=2\sqrt{3}
	\end{aligned}$

	La distance entre C et $P$, c'est-à-dire la distance de C à son projeté orthogonal sur $P$ est donc bien égale à $2\sqrt{3}$.



\end{enumerate}
\medskip

\end{document}
